\section{CONCLUSION AND BOUNDARIES}\label{sec:conclusion}

The paper set out to answer four questions.

\emph{What can the finding record tell us?} A finding history predicts the
future finding record. It fixes a floor under the underlying condition rate
in every history group and, if follow-up can only raise detection, a ceiling
on the underlying-condition risk ratio: the observed finding ratio. Where the
observed ratio is below one, the underlying association is negative.

\emph{What can it only bound?} The underlying-condition risk ratio itself.
Without information about detection, the record places it anywhere between
$q_1$ and $1/q_0$, an interval that at ordinary finding rates contains one and
runs well below it. A longer record does not narrow the interval, because
every added year is another observation of the product $q = rd$.

\emph{Why is follow-up different from discovery?} A prior finding is routing
information: it says where to return. The follow-up procedures, not the prior
finding, produce the evidence about whether the known condition remains.
Finding histories can contain genuine information about both recurrence and
other problems. But when the history also changes future observation, the
record alone cannot determine how much of that relationship belongs to the
underlying condition and how much belongs to the observation process.
Follow-up uses history to revisit known conditions. Discovery requires
observation whose measurement rule is not itself distorted by the history
being evaluated, or a defensible model of that distortion. In
Clearinghouse data, prior findings predict recurrence and materially more
newly recorded findings, and the record does not identify why.

\emph{What additional audit information would let us learn more?} Anything
that pins down the follow-up detection multiplier $M$ for the finding class
in question: reference observation under a common measurement rule, which
recovers the underlying ratio and hence $M$ when the reference procedure's
sensitivity is common to both history groups; a high-sensitivity reference
procedure that anchors the condition rate; reference adjudication accurate
enough to stand in for latent status; two procedures with independent
failures; quasi-random changes in attention; or, at minimum, a defended
ceiling $\Mbar$, under which the underlying ratio is at least $\Robs/\Mbar$.
Each is an addition to audit design, not a statistical method applied to the
record.

Four boundaries travel with these conclusions. Repeat findings are not
artifacts; a repeat finding in these data is a real reported deficiency the
auditor identified as a recurrence. A newly recorded finding is a problem
recorded for the first time in this program and year, a statement about the
record and not about when the problem arose. The predictive results are
constructions showing what can happen when a score influences observation;
whether a given deployed model has this property is an empirical question
that in most cases nobody has checked. And no result in Section~\ref{sec:fac}
identifies the mechanism behind the associations it reports. Nothing here
recommends a particular protocol for a particular office.

The missing quantity is detection. Once an organization measures, bounds, or
standardizes it, the finding record becomes much more informative about the
underlying condition.
